\documentclass[a3paper,landscape,9pt]{article}

\usepackage[utf8]{inputenc}
\usepackage[T1]{fontenc}
\usepackage[french]{babel}
\usepackage{lmodern}
\usepackage[a3paper,landscape,margin=0.9cm]{geometry}
\usepackage{xcolor}
\usepackage{tikz}
\usepackage{pgfplots}
\usepackage{booktabs}
\usepackage{array}
\usepackage{enumitem}
\usepackage{ragged2e}
\usepackage{microtype}

\pgfplotsset{compat=1.18}
\pagestyle{empty}
\renewcommand{\familydefault}{\sfdefault}
\setlength{\parindent}{0pt}
\setlength{\fboxsep}{5pt}
\setlength{\fboxrule}{0.55pt}
\setlist[itemize]{leftmargin=*, itemsep=1pt, topsep=1pt}
\sloppy
\emergencystretch=3em

\definecolor{navy}{HTML}{0B2F4A}
\definecolor{blue}{HTML}{2F80ED}
\definecolor{darkblue}{HTML}{145DA0}
\definecolor{pale}{HTML}{F7FBFE}
\definecolor{cyan}{HTML}{EAF4FC}
\definecolor{line}{HTML}{BFD3E2}
\definecolor{green}{HTML}{1F9D7A}
\definecolor{orange}{HTML}{F2994A}
\definecolor{textgray}{HTML}{26323D}
\definecolor{muted}{HTML}{52616B}

% Cadre à hauteur contrôlée.
% La largeur interne retire l'épaisseur de la bordure et l'espace du fcolorbox :
% cela évite que le texte ou la bordure dépasse de la colonne.
\newcommand{\posterbox}[3]{%
  \fcolorbox{line}{pale}{%
    \begin{minipage}[t][#1][t]{\dimexpr\linewidth-2\fboxsep-2\fboxrule\relax}
      \RaggedRight
      {\large\bfseries\color{navy} #2}\par
      \vspace{0.10cm}
      {\footnotesize #3}
    \end{minipage}%
  }%
}

\newcommand{\metric}[4]{%
  \begin{minipage}[t]{#1}
    \fcolorbox{#4}{white}{%
      \begin{minipage}[c][1.10cm][c]{\dimexpr\linewidth-2\fboxsep-2\fboxrule\relax}
        \centering
        {\Large\bfseries\color{#4} #2}\par
        {\scriptsize\bfseries #3}
      \end{minipage}%
    }%
  \end{minipage}%
}

\newcommand{\stepbox}[3]{%
  \begin{minipage}[t]{0.235\linewidth}
    \fcolorbox{line}{white}{%
      \begin{minipage}[c][0.82cm][c]{\dimexpr\linewidth-2\fboxsep-2\fboxrule\relax}
        {\bfseries\color{blue} #1}\quad{\bfseries #2}\par
        {\scriptsize\color{muted} #3}
      \end{minipage}%
    }%
  \end{minipage}%
}

\begin{document}

{\Huge\bfseries\color{navy} Projet Fashion-MNIST : classification et clustering}\par
\vspace{0.08cm}
{\large LU3IN0226 - IA \& Data science \hfill Noms : \textbf{Mehdi Qostali et Hakim Ziane}}\par
\vspace{0.05cm}

\vspace{0.22cm}

\fcolorbox{blue}{cyan}{%
  \begin{minipage}[c][2.85cm][c]{\dimexpr\linewidth-2\fboxsep-2\fboxrule\relax}
    \RaggedRight
    {\Large\bfseries\color{navy} Lecture rapide}\par
    \vspace{0.14cm}
    {\large\bfseries La forêt aléatoire est le meilleur modèle supervisé testé : 0,8535 en binaire et 0,8700 en multi-classe.}\par
    \vspace{0.12cm}
    {\small Les erreurs viennent surtout des vêtements du haut : \textit{Shirt}, \textit{T-shirt/top}, \textit{Pullover} et \textit{Coat}.}\par
    \vspace{0.22cm}
    \stepbox{1}{Données}{images + labels}
    \hfill
    \stepbox{2}{Préparation}{normalisation}
    \hfill
    \stepbox{3}{Modèles}{supervisé + clusters}
    \hfill
    \stepbox{4}{Évaluation}{accuracy, ARI, NMI}
  \end{minipage}
}

\vspace{0.20cm}

\begin{center}
\metric{0.23\linewidth}{0,8535}{Accuracy test binaire}{green}
\hfill
\metric{0.23\linewidth}{0,8700}{Accuracy test multi-classe}{blue}
\hfill
\metric{0.23\linewidth}{0,611}{NMI hiérarchique}{orange}
\hfill
\metric{0.23\linewidth}{86,6\,\%}{Variance gardée par l'ACP}{darkblue}
\end{center}

\vspace{0.25cm}

\begin{minipage}[t]{0.31\linewidth}
\posterbox{3.10cm}{Données et objectif}{
Fashion-MNIST contient 10 classes : T-shirt/top, Trouser, Pullover, Dress, Coat, Sandal, Shirt, Sneaker, Bag et Ankle boot.

\vspace{0.08cm}
Le but est de comparer plusieurs modèles supervisés, puis de voir si le clustering retrouve des groupes cohérents sans utiliser les labels pendant l'apprentissage.
}

\vspace{0.18cm}

\posterbox{3.65cm}{Préparation}{
Chaque image est représentée par 784 pixels et les valeurs sont normalisées dans $[0,1]$.

\vspace{0.08cm}
Les validations croisées sont faites sur des échantillons stratifiés : chaque classe reste représentée dans les mêmes proportions.

\vspace{0.08cm}
Le test final utilise \texttt{fashion-mnist\_test.csv}.
}

\vspace{0.18cm}

\posterbox{2.95cm}{Lecture des matrices}{
Ligne = vraie classe. Colonne = classe prédite.

\vspace{0.08cm}
Bleu clair = peu d'exemples. Bleu foncé = beaucoup d'exemples.

\vspace{0.08cm}
La diagonale indique les bonnes prédictions. Les autres cases indiquent les erreurs.
}

\vspace{0.18cm}

\posterbox{3.00cm}{À retenir}{
Les classes les plus simples sont celles dont la silhouette est très différente, par exemple les sacs, les pantalons et plusieurs chaussures.

\vspace{0.08cm}
Les vêtements du haut se ressemblent davantage, ce qui explique la plupart des confusions.
}
\end{minipage}
\hfill
\begin{minipage}[t]{0.34\linewidth}
\posterbox{3.15cm}{Résultats supervisés}{
\textbf{Binaire.} La tâche T-shirt/top vs Shirt est difficile, car les formes sont proches. La forêt aléatoire atteint 0,8535 sur le test.

\vspace{0.08cm}
\textbf{Multi-classe.} Sur les 10 classes, la forêt aléatoire atteint 0,8700 après entraînement sur 20\,000 exemples.
}

\vspace{0.18cm}

\posterbox{6.05cm}{Comparaison des classifieurs}{
\begin{center}
\begin{tikzpicture}
\begin{axis}[
    width=0.96\linewidth,
    height=4.95cm,
    ybar,
    ymin=0,
    ymax=1,
    ylabel={Accuracy CV},
    symbolic x coords={Baseline,Arbre,Perceptron,kNN,SGD log.,Forêt},
    xtick=data,
    x tick label style={rotate=30,anchor=east,font=\scriptsize},
    legend style={at={(0.02,0.98)},anchor=north west,draw=none,font=\scriptsize},
    bar width=5pt,
    grid=major,
    grid style={line!55},
]
\addplot[fill=blue] coordinates {
    (Baseline,0.100) (Arbre,0.721) (Perceptron,0.776) (kNN,0.791) (SGD log.,0.802) (Forêt,0.834)
};
\addplot[fill=green] coordinates {
    (Baseline,0.500) (Arbre,0.790) (Perceptron,0.801) (kNN,0.828) (Forêt,0.842)
};
\legend{Multi-classe,Binaire}
\end{axis}
\end{tikzpicture}
\end{center}
La forêt aléatoire donne les meilleurs scores parmi les modèles testés.
}

\vspace{0.18cm}

\posterbox{4.65cm}{Erreurs les plus fréquentes}{
\begin{tabular}{@{}p{0.66\linewidth}r@{}}
\toprule
Vraie classe $\rightarrow$ classe prédite & Nb \\
\midrule
Shirt $\rightarrow$ T-shirt/top & 179 \\
Pullover $\rightarrow$ Coat & 113 \\
Shirt $\rightarrow$ Pullover & 112 \\
Shirt $\rightarrow$ Coat & 90 \\
T-shirt/top $\rightarrow$ Shirt & 85 \\
\bottomrule
\end{tabular}

\vspace{0.10cm}
La classe \textit{Shirt} est la plus fragile : son rappel est autour de 0,566.
}

\vspace{0.18cm}

\posterbox{2.75cm}{Interprétation}{
La forêt aléatoire sépare bien les classes très différentes, mais elle hésite quand plusieurs vêtements ont une forme globale proche en $28 \times 28$ pixels.
}
\end{minipage}
\hfill
\begin{minipage}[t]{0.31\linewidth}
\posterbox{3.30cm}{Non-supervisé}{
Avant le clustering, une ACP réduit la dimension.

\vspace{0.08cm}
Pour k-means, 50 composantes principales gardent environ 86,6\,\% de la variance.

\vspace{0.08cm}
Pour le hiérarchique, on utilise 500 exemples et 20 composantes pour limiter le coût.
}

\vspace{0.18cm}

\posterbox{6.05cm}{Scores de clustering}{
\begin{center}
\begin{tikzpicture}
\begin{axis}[
    width=0.96\linewidth,
    height=4.95cm,
    xmin=4,
    xmax=21,
    ymin=0,
    ymax=0.62,
    xlabel={Nombre de clusters},
    ylabel={Score},
    xtick={5,10,15,20},
    grid=major,
    grid style={line!55},
    legend style={at={(0.02,0.98)},anchor=north west,draw=none,font=\scriptsize},
]
\addplot[color=green,mark=*] coordinates {(5,0.199) (10,0.199) (15,0.174) (20,0.162)};
\addplot[color=blue,mark=*] coordinates {(5,0.284) (10,0.385) (15,0.373) (20,0.326)};
\addplot[color=orange,mark=*] coordinates {(5,0.423) (10,0.536) (15,0.547) (20,0.538)};
\addplot[color=muted,dashed] coordinates {(10,0) (10,0.62)};
\legend{Silhouette,ARI,NMI,$k=10$}
\end{axis}
\end{tikzpicture}
\end{center}
$k=10$ est retenu pour correspondre aux 10 classes Fashion-MNIST.
}

\vspace{0.18cm}

\posterbox{3.45cm}{Résultats de clustering}{
\textbf{k-means :} ARI $\simeq 0,385$ et NMI $\simeq 0,536$ avec $k=10$.

\vspace{0.08cm}
\textbf{Hiérarchique :} ARI $\simeq 0,430$ et NMI $\simeq 0,611$ sur 500 exemples.

\vspace{0.08cm}
Les groupes sont cohérents visuellement, mais ne recouvrent pas parfaitement les 10 labels.
}

\vspace{0.18cm}

\posterbox{3.30cm}{Bilan et pistes}{
Le supervisé reste plus performant pour prédire les labels.

\vspace{0.08cm}
Pour améliorer : entraîner sur toute la base, optimiser les hyper-paramètres, tester un modèle convolutionnel et analyser séparément Shirt, Pullover, Coat et T-shirt/top.
}
\end{minipage}

\vfill

\end{document}
